\begingroup
\linespread{1.22}\selectfont
\setlength{\parskip}{.32em}
\setcounter{figure}{0}
\renewcommand{\thefigure}{概.\arabic{figure}}
\chapter*{概论：从业务语义看分布式系统}
\phantomsection\label{front:overview}
\addcontentsline{toc}{chapter}{概论：从业务语义看分布式系统}

\begin{center}
  \includegraphics[width=\linewidth,height=.35\textheight,keepaspectratio]{assets/chapter-art/overview.png}
  \par\vspace{4pt}
  {\scriptsize\color{Muted}取意于彼得·勃鲁盖尔《雪中猎人》\par}
  \BookEpigraph{overview}
\end{center}

一项在线购买服务，对用户呈现的是一个连续过程：确认商品与金额，提交购买，查看付款和发货，必要时申请退款。支撑这个过程的系统却分布在许多位置。入口接收请求，数据库保存订单，支付平台改变资金，消息系统传递已经发生的事实，仓库安排发运。用户稍后打开页面，看到的状态还可能来自另一份查询数据和缓存。

这些位置共同提供服务，又分别拥有自己的处理速度、提交时刻和故障。购买进程退出时，支付平台可以继续完成收款。仓库暂停以后，用户可以先查到付款成功。某个地域恢复较早的备份，其他地域的交易与外部平台的账本又可能保留较新的历史。一项业务的含义，需要在这些分散的位置和不同的时刻之间保持清楚。

本书以幂等性为线索，研究这些关系怎样成立。一次确认可能产生多次请求，一项工作可能经由多个执行者接力，已经发生的效果也可能早于任何一方取得确定答复。沿着同一项购买，我们会从业务身份、并发仲裁与本地事务，走到跨系统协作、历史保存和完整架构，逐步说明分布式系统怎样把局部活动组织成可以理解、可以接续的服务。

\section*{把一项服务放进完整系统}
\pdfbookmark[1]{把一项服务放进完整系统}{overview-panorama}

系统分布到多台机器，首先有实际的能力需求。访问量增长以后，计算、内存和磁盘需要分担。单个节点退出时，其他位置要保存能够继续使用的事实。用户、设备与仓库分布在不同地域，服务需要安排数据和处理的位置。支付机构等外部参与者还拥有独立的账本、授权和运行规则。分布带来了容量、容错、局部性和隔离，也使一项服务开始依赖多个能够独立变化的部分。

可以先把购买系统看成几条相互联系的路径。在线业务从用户确认进入受理，再根据商品和资源条件取得库存资格，向平台付款，形成订单并推进履约。查询路径回答原购买结果、当前履约与退款进度等问题。恢复路径在前台连接结束以后继续发现未完成的工作，从平台、消息和其他保存位置取得证据，再推进相应责任。图~\ref{fig:overview-panorama} 把这些路径及其支撑设施放在一起。

\begin{figure}[htbp]
\centering
\begin{tikzpicture}[
  box/.style={draw=Ink!50,fill=Paper,rounded corners=2pt,
    text width=3.5cm,minimum height=1.12cm,align=center,inner sep=5pt,
    font=\sffamily\fontsize{9}{12}\selectfont},
  authority/.style={box,draw=Accent!75,fill=Accent!7},
  support/.style={box,fill=Ink!4,minimum height=1.05cm},
  arr/.style={-{Stealth[length=1.8mm]},draw=Ink,line width=.7pt},
  evidence/.style={arr,dashed,draw=Muted},
  tag/.style={font=\sffamily\fontsize{8}{10}\selectfont,
    text=Muted,fill=white,inner sep=2pt,align=center}
]
\node[box] (entry) at (0,0) {购买与查询入口\\用户意图、身份与授权};
\node[authority] (purchase) at (5.0,0) {购买权威\\受理、原结果与未决责任};
\node[authority] (platform) at (10.0,0) {支付平台\\收款、退款与外部账本};
\node[box] (query) at (5.0,2.45) {查询投影与缓存\\原结果与当前进度的观察};
\node[support] (control) at (10.0,2.45) {控制面\\归属、成员与执行资格};
\node[authority] (inventory) at (0,-2.65) {库存分区\\资源占用与联合决定};
\node[box] (message) at (5.0,-2.65) {消息域\\事件保存、传输与位置};
\node[authority] (fulfill) at (10.0,-2.65) {履约与发运\\本地接受、外部发货};
\draw[arr] (entry)--node[tag,above=8mm]{受理与恢复}(purchase);
\draw[arr] (purchase)--node[tag,above=8mm]{原付款／退款身份}(platform);
\draw[arr] (purchase.south west)--node[tag,sloped,above]{资源资格}(inventory.north east);
\draw[arr] (purchase)--node[tag,right]{已保存事件}(message);
\draw[arr] (message)--node[tag,above=8mm]{交付与接受}(fulfill);
\draw[evidence] (purchase)--node[tag,right]{派生观察}(query);
\draw[evidence] (entry.north) |- node[tag,pos=.75,above]{读取观察} (query.west);
\draw[evidence] (control.south west)--node[tag,sloped,above]{位置与资格}(purchase.north east);
\draw[evidence] (fulfill.south) -- ++(0,-.8) -|
  node[tag,pos=.25,below]{持久交接回执} (purchase.south east);
\node[support] at (0,-5.15) {完整图像与归档\\保存恢复材料和来源};
\node[support] at (5.0,-5.15) {任务与持久执行\\调度工作者、保存执行位置};
\node[support] at (10.0,-5.15) {部署、容量与运维\\故障隔离和恢复预算};
\end{tikzpicture}

\caption{一项服务由业务、观察与恢复路径共同构成，各类设施承担不同的保存和执行职责。}
\label{fig:overview-panorama}
\end{figure}

图中的购买权威，是有权决定购买是否受理、原身份绑定什么输入、原结果是什么的保存位置。权威描述一种职责：它可以由数据库中的约束和事务承担，也可以由一组副本共同执行的状态机承担。状态机根据当前状态和输入命令产生下一状态，同一组副本按确定的提交顺序应用命令，保持相同的业务关系。履约数据库、库存分区和支付平台也分别决定自己掌管的事实。

这些权威之间有业务关联。购买记录需要关联实际收款，付款成功要形成可以传播的事件，消费者接受履约以后又要留下可查询的接受事实。每一方的记录都能解释自己已经确认了什么，保存下来的关联使恢复者能把它们接回同一项购买。一份局部记录是否可信，取决于它由谁产生、保存过什么决定，以及经历故障以后保留到了哪里。

查询数据与缓存处在另一条路径。页面可以把付款、履约和退款组合成便于读取的视图，再把常用观察放到靠近用户的位置。这类从原事实计算出来的状态称为派生数据。缓存丢失时，可以依据可靠来源重建。缓存落后时，页面所见可能晚于实际提交，系统据其用途决定怎样标明观察时间，以及哪类判断必须回到当前权威。

计算组件负责执行工作。入口进程、支付工作者、事件发布器、消费者和对账作业可以各自启动、暂停和更换。工作者的内存可以保存一次执行的局部进度，跨进程的恢复则从持久责任与原操作关联进入。这样，一笔付款等待数小时，承担它的某个进程可以只运行数分钟。

此外，系统要知道逻辑分区当前放在哪里、哪个副本或工作者有资格写入、迁移何时完成。这类管理位置、成员和资格的系统称为控制面。承载业务读写与执行的路径则组成数据面。控制面更新归属，实际接收写入的资源核验资格，入口缓存减少每次访问对配置服务的依赖。配置暂时不可用时，哪些既有路径能够继续，由这些职责的安排决定。

存储把大规模日志、备份、图像和对账材料保存下来。这里的图像指某个明确位置上的完整系统状态，并包含恢复它所需的关联信息。存储服务提供字节和版本的保存能力，发布规则再确定哪一组字节构成有效状态。多地域部署安排这些保存与执行位置的故障关系，容量和运维使在线请求、后台恢复与数据搬运获得实际资源。

\section*{业务含义、执行尝试与确定结果}
\pdfbookmark[1]{业务含义、执行尝试与确定结果}{overview-semantics}

全景中的第一项区别来自用户意图。Alice 点击一次确认，连接超时以后再次提交，两份报文可能属于同一项购买。她稍后确实想再买一份，即使商品、数量和金额相同，也应形成另一项购买。只比较报文内容会丢失这两种情况之间的区别。系统需要在调用方决定意图时保存身份，并把同一意图的后续尝试接到这个身份上。

本书把一次调用、一次消息投递或一次任务运行称为执行尝试。尝试可以在任何一步结束，也可以与另一个尝试同时存在。业务身份表示它们共同处理的工作。第一部分会具体建立账户、动作、会话和客户端键组成的完整身份，解释各分量由谁保存、怎样传播、何时允许新受理。相同身份下的原输入一旦形成提交绑定，后续恢复就据此比较。同一身份改成另一金额时，服务明确报告冲突，已经成立的购买继续保留原来的含义。

身份还要进入真正能够限制业务变化的位置。两个并发请求都先查到空记录，随后各自扣款，说明一次读取尚未决定谁能首次执行。数据库唯一约束可以让相同身份在同一受理位置竞争，业务事务再把身份、资金变化、订单和原结果共同保存。所谓原子提交，表示这些相关变化作为一组对外成立，或者在提交前失败时共同撤销。账户余额、订单归属与实际金额又由各自的业务规则检查。

可以用一份独立的本地购买看清结果。Alice 初始有 10000 分，购买金额为 3000 分，成功以后剩 7000 分。这次购买的原结果保存当次扣款及相应编号。随后另一项合法支出扣去 1000 分，当前余额变成 6000 分。恢复原购买时，原结果中的当次余额继续为 7000 分，账户查询则报告当前的 6000 分。图~\ref{fig:overview-identity} 将尝试、仲裁与这两种观察分别画出。

\begin{figure}[htbp]
\centering
\begin{tikzpicture}[
  box/.style={draw=Ink!45,fill=Paper,rounded corners=2pt,
    text width=3.3cm,minimum height=1.05cm,align=center,inner sep=5pt,
    font=\sffamily\fontsize{9}{12}\selectfont},
  original/.style={box,draw=Accent!75,fill=Accent!7},
  arr/.style={-{Stealth[length=1.8mm]},draw=Ink,line width=.7pt},
  note/.style={font=\sffamily\fontsize{8}{11}\selectfont,
    text=Muted,align=center,text width=3.4cm}
]
\node[box] (a1) at (0,0) {尝试一\\提交以后响应丢失};
\node[box] (a2) at (0,-1.65) {尝试二\\携带原身份与原输入};
\node[box] (a3) at (0,-3.3) {后续恢复\\更换进程或连接};
\node[original,text width=3.8cm,minimum height=3.0cm] (decision) at (5,-1.65)
  {同一购买的仲裁位置\\[3pt]原身份绑定原输入\\共同保存业务效果\\保存实际编号与原结果};
\node[original] (result) at (10,0) {原购买结果\\扣款后余额 7000 分};
\node[box] (current) at (10,-3.3) {当前账户查询\\另支出后余额 6000 分};
\draw[arr] (a1.east)--(decision.north west);
\draw[arr] (a2)--(decision);
\draw[arr] (a3.east)--(decision.south west);
\draw[arr] (decision.north east)--(result.west);
\draw[arr] (decision.south east)--(current.west);
\node[note] at (0,1.35) {用户只确认一次购买\\多次通信尝试继续引用它};
\node[note] at (5,-4.25) {独立初态 10000 分\\本次购买扣款 3000 分};
\node[box,minimum height=.75cm] (later) at (10,-1.65)
  {另一合法支出 1000 分};
\draw[arr,draw=Muted] (later.south)--(current.north);
\end{tikzpicture}

\caption{同一购买的多次尝试沿原身份汇合，原结果与后续变化形成的当前状态分别保存和查询。}
\label{fig:overview-identity}
\end{figure}

相同的区分也适用于外部付款。平台已经收款，是平台的事实；购买服务取得并保存这份证据，是另一次提交；用户收到响应，又是后续的通信结果。付款成功、履约接受和包裹实际发出还各有自己的时刻。接口需要说明它报告的是哪一层事实，恢复者才能知道下一步该查询、等待还是继续推进。

幂等性在数学上描述一种重复性质：对同一状态应用某个操作，再应用它，所得状态与只应用一次相同。业务中还会有其他操作穿插，返回值也承载历史含义。例如把余额设置成 7000，紧接着再设一次会保持 7000。中间若发生一笔合法的 1000 分支出，旧设置再次到达却会把 6000 改回 7000。连续重复的性质与交错历史中的正确行为，需要分别分析。

交换性讨论两个操作互换次序是否得到相同的相关状态和结果。单调性讨论状态是否沿选定的信息顺序增长。合并收敛讨论不同副本在取得相同信息以后能否达到共同状态。业务不变量则规定每个允许状态都应保持的关系，例如余额非负、库存守恒、退款占用不超过原收款。选择适合的状态表达可以减少协调或历史记录，资金、授权与资源限制继续由对应规则约束。

\section*{从局部提交走向跨域协作}
\pdfbookmark[1]{从局部提交走向跨域协作}{overview-cooperation}

购买和扣款在同一个数据库时，可以通过一项事务共同提交。接入独立支付平台以后，本地受理和平台收款分属两个保存位置。购买服务可能已经收到平台成功，尚未把结果写回就退出；也可能请求尚未到达平台，先失去连接。恢复者读到的本地处理中状态相同，平台上的正确后续却取决于先前请求是否已经发生效果。

未知结果由此成为一种需要持续管理的业务状态。它表达当前保存位置缺少确定结论。购买服务先保存已经受理的原输入、平台身份和后续责任，恢复者就能沿同一关联继续查询、接收回调或对账。平台若在约定范围与期限内对原键和输入进行仲裁，重发可以接回原收款。平台只提供某个时点的查找结果时，一次未查到之后，较早发出的请求还可能到达，发送决策就要依据这项能力安排。

持久工作流把接续过程保存下来。工作流可以记录业务状态、下一步动作、查询时间和当前执行资格。进程退出以后，扫描器从这些状态重新发现工作。接管者取得新的工作代次，资源端把它与实际写入共同检查，迟到执行者就无法凭过时资格覆盖当前决定。已经离开旧进程的支付请求，则由原平台身份与后续证据连接，接管记录同时说明谁能继续安排本地工作。

这种安排使授权、执行和结果各有依据。当前工作者可以换人，同一购买继续使用原身份。可信的旧响应可能带来已经发生的事实，是否接受其中证据，要核对原对象、金额、币种与来源。执行代次约束谁能提交相应更新，业务版本解释同一责任推进到了哪一步。第八章会具体展开租约、接管和资源端检查怎样配合。

付款成功之后，还要把事实交给履约。如果先改数据库、再发送消息，服务可能在两步之间退出。把待发事件与业务结果放进同一次提交，可以先保存传播责任。这份事务内的待发事件记录通常称为 Outbox。发布器稍后把它交给消息系统，确认丢失时可以继续传送原事件。消息系统中的投递标记或日志位置可能变化，业务事件身份随原记录保持。

消费者也有自己的提交边界。它可以在同一事务内保存接收记录、初次履约的业务守卫、实际履约编号和发运意图。接收记录通常称为 Inbox，用来识别某事件是否已被本地处理。业务守卫再限制诸如同一订单的初次履约这类业务关系，即使两个不同事件都声称推进它，也会在相同的业务约束处汇合。消费确认、原结果与后续交接分别结束相应责任。

消息传播存在存储和保留条件。消息域接受一份事件以后，承担约定范围内的保存与交付；消费者处理完毕以后，用确认或位移记录推进传输位置。保留时间有限时，上游还需要取得可信的持久交接事实，才能结束原事件的补发责任。这样，发送者和接收者各自能解释已经交给谁、哪段工作已经有人保存。

另一些业务需要多个资源先形成共同决定。例如组合商品同时依赖两个库存分区，付款资格就要建立在约定的联合预留之上。参与者保存准备占用，协调者保存唯一的提交或撤销决定，恢复者传播该决定并确认必需参与者完成。这里的分布式事务解决资源之间确实需要共同判断的部分，外部付款和发运再沿自己的责任与证据继续。

需要退款时，系统创建关联原收款的新动作。收款 3000 分，已经成功退回 1000 分，另有 1500 分退款结果未知，后续最多再占用 500 分。未知退款继续占用相应额度，取得成功或足够最终拒绝的证据以后再转换或释放。补偿动作因此有自己的身份、授权、结果和恢复路径，原购买成功与当前净资金也能分别解释。

这些协作共有一种结构：发送前保存需要完成的工作及其身份，接收者在能控制本地效果的位置仲裁，确定结果沿原关联返回。一次通信可以结束，持久责任继续存在。每个跨域边界都按实际能力说明已知事实和允许动作，整项业务便能由多个局部过程接续完成。

\section*{系统的过去怎样延续}
\pdfbookmark[1]{系统的过去怎样延续}{overview-history}

前面的恢复都依赖系统还记得已经做过什么。一次成功结果保存在主库，主库退出以后接班副本如果缺少这项记录，旧身份就可能再次遇到空位置。复制要保存与已确认业务相匹配的历史，接班者才能沿同一身份恢复。Paxos、Raft 等共识协议使一组副本在相应故障模型中确定共同决定，业务状态机再按提交顺序把输入、结果、会话和未决责任应用到自己的状态。

复制与事务在这里承担不同关系。事务安排哪些修改共同成立，复制安排已确认状态怎样跨节点故障保存。复制日志中的一条命令可以在状态机中完成一组关联修改，数据库也可以在内部用复制保存事务决定。跨多个组的资源变化，则另外安排共同决定和恢复。观察其中一层的成功，需要知道它在整个提交路径中的位置。

快照改变历史的表示。长日志可以压缩成某个已应用位置的状态图像，再从后续日志继续。原身份、原结果、会话资格、未决工作和编号进度属于恢复判断的一部分。完整图像连同应用边界共同保存，接班者才能知道自己从哪一个确定位置继续。只保存账户金额，恢复以后就可能缺少解释这笔金额为何变化的关系。

回收又会改变可以取得的细节。服务可以约定旧提交入口的有效期限，关闭会话以后按规则减少逐项结果。关闭决定与必要拒绝依据继续保存，较老的请求才有明确处置。已经受理的付款责任由服务端接续，客户端恢复入口的结束与业务责任的结束按各自条件处理。平台键保护期、在途请求和外部可查证据也参与决定哪些动作仍有依据。

业务表达能够改变保存历史的成本。标签按集合并入，重复传输同一成员可以自然汇合；带来源的新增与删除标记又能保留并发差异。相应结构称为无冲突复制数据类型，即 CRDT，它在规定的状态和传输条件下安排副本合并。资金扣款、退款占额与初次发运涉及另外的业务约束。两份副本即使最后收敛，也可能各自先接受了合计超出资源的工作。

在线迁移改变承载位置。逻辑分区表示业务事实的稳定归属，副本组表示它当前由哪些机器保存和执行。迁移把完整事实与未决责任交给新位置，再完成执行资格的切换，已有购买沿原关联继续。控制面告诉入口去哪里，源与目标的资源端共同约束哪个位置可以接受新的写入。

灾难回退进一步改变历史本身。本地从较早备份重建，外部平台的收款与履约库的发货可能已经向前。恢复程序先找回原身份到原付款的绑定，核对较新的外部事实，再按已经恢复完整的范围开放入口。缺少可靠关联的旧范围，可以通过持久关闭规则限制新的提交。新的编号代次安排新工作的编号来源，旧工作从故障前的原身份与原关联恢复。

图~\ref{fig:overview-history} 将这些变化分别放到保存与恢复路径上。副本、快照和迁移处理怎样把既有事实带到后继位置，回收安排旧入口与剩余证据，旧备份恢复则重新对齐已经分散的业务历史。

\begin{figure}[htbp]
\centering
\begin{tikzpicture}[
  box/.style={draw=Ink!45,fill=Paper,rounded corners=2pt,
    text width=3.5cm,minimum height=1.1cm,align=center,inner sep=5pt,
    font=\sffamily\fontsize{9}{12}\selectfont},
  original/.style={box,draw=Accent!75,fill=Accent!7},
  arr/.style={-{Stealth[length=1.8mm]},draw=Ink,line width=.7pt},
  tag/.style={font=\sffamily\fontsize{8}{10}\selectfont,
    text=Muted,fill=white,inner sep=2pt,align=center}
]
\node[original,minimum height=2.0cm] (facts) at (0,0)
  {已受理业务的完整状态\\身份、输入与原结果\\未决责任、会话与编号};
\node[box] (replica) at (5,1.65) {副本接班\\保留已确认的提交历史};
\node[box] (snapshot) at (5,0) {快照与压缩\\保存完整状态及应用位置};
\node[box] (migration) at (5,-1.65) {在线迁移\\原事实与执行资格交接};
\node[original,minimum height=2.0cm] (continue) at (10,0)
  {后继位置继续原业务\\取得原关联与结果\\接续尚未完成的工作};
\draw[arr] (facts.north east)--(replica.west);
\draw[arr] (facts)--(snapshot);
\draw[arr] (facts.south east)--(migration.west);
\draw[arr] (replica.east)--(continue.north west);
\draw[arr] (snapshot)--(continue);
\draw[arr] (migration.east)--(continue.south west);
\node[box] (gc) at (5,-3.5)
  {回收逐项历史\\关闭旧准入，保留拒绝依据\\未决责任继续保存};
\draw[arr] (facts.south) |- (gc.west);
\draw[arr] (gc.east) -| (continue.south);
\node[box] (old) at (0,-6.0)
  {旧备份重新出现\\外部可能保留较新效果};
\node[original] (reconcile) at (5,-6.0)
  {恢复原关联与外部证据\\或持久关闭歧义范围};
\node[box] (admission) at (10,-6.0)
  {按已核准范围开放\\旧工作沿原身份恢复};
\draw[arr] (old)--(reconcile);
\draw[arr] (reconcile)--(admission);
\node[tag] at (5,-4.75) {灾难回退按各保存域的实际历史重新安排恢复与准入};
\end{tikzpicture}

\caption{保存表示和承载位置可以变化，后继处理依据原业务关联、剩余证据与明确的准入规则继续。}
\label{fig:overview-history}
\end{figure}

\section*{基础设施怎样承担这些关系}
\pdfbookmark[1]{基础设施怎样承担这些关系}{overview-infrastructure}

完整架构要为这些关系安排具体承担位置。首先是数据分区：哪些账户、购买、会话、结果和责任由同一个逻辑分区保存，哪些资源位于其他分区。合理分区可以让常见修改在局部完成，也能分担总容量。共享余额、热门库存或编号分配器会形成热点，改变机器数量以后，这些共享点的竞争关系也会继续影响吞吐。

复制组承载每个分区的决定。Raft 通过日志、任期与多数提交组织一组副本，链式复制则按头部接收、沿链传播和尾部确认等规则安排状态与读取。复制顺序怎样进入业务应用，读者取得的是哪个确定位置，成员变化时怎样交接，都直接影响对外观察。TiKV 把有序键空间划成 Region，每个 Region 由一组副本保存，跨这些范围的事务再由上层机制组织。

控制面保存成员、配置、归属和资格。ZooKeeper 用节点、版本与会话组织协调，etcd 提供带修订号的键值状态、事务与租约，Consul 也有自己的服务发现、会话和读取规则。应用从这些接口取得管理事实，缓存通知以后还要能重建完整观察。资格最终在实际修改资源的位置核验，配置变化才能对应到谁可以执行什么。

消息设施安排通信与时间的解耦。RabbitMQ 把消息经交换机路由到队列，由确认推进相应交付责任，适合围绕待处理工作组织消费。Kafka 把记录追加到分区日志，不同消费组分别保存读取位置，可以在保留期内独立重读。两者各有复制、确认、并发与保留规则。购买事件、业务接收和实际履约则沿保存的身份与本地事务对应。

Redis 可以为查询提供快速的键值观察。理解这种位置，需要深入命令怎样执行、修改何时进入持久化与复制路径、缓存过期与淘汰怎样工作。RDB 保存快照，AOF 保存用于重放的修改历史，副本切换和集群迁槽又会改变位置与可见状态。把 Redis 用作派生缓存时，故障后从权威重建；选择让某项业务决定由它承担时，就要把相应的共同修改、持久性和恢复条件一起安排。

分布式存储承担大规模字节与完整版本的保存。Google 文件系统 GFS 把文件内容划分成块，由元数据服务管理位置，客户端直接向数据节点传输内容。Hadoop 分布式文件系统 HDFS 也区分命名与数据路径，并有自己的写入流水线和恢复规则。对象存储通过对象与版本提供独立的持久容量。日志、归档和状态图像使用它们时，要按实际接口解释何时可读、怎样重试、哪些字节属于同一份有效材料。

图像与批量输出通常先形成候选内容，再通过一项完整清单发布有效集合。清单记录来源、版本、块顺序和校验信息，恢复者固定清单以后取得完整字节。计算过程可以重做，发布位置负责把成功输出接到逻辑工作上。这样，任务运行、文件生成与结果正式可见在架构中各有清楚的边界。

分布式数据库把其中一些关系组织成通用数据接口。SQL 是描述关系数据查询与修改的语言，数据库在它的下面安排索引、复制、隔离和事务。Spanner 将数据复制、跨组提交和有界时间不确定性连接起来；CockroachDB 用范围副本、写意图和事务记录组织提交与恢复；TiDB 在 TiKV 的分区结构上执行关系查询与事务。应用把完整业务身份、资源约束和待办映射进所选数据库的事务，依其错误与重试规则继续。

FaRM 展示另一种实现位置。它把分布式内存、远程直接内存访问技术和事务日志连接起来。远程直接内存访问简称 RDMA，允许网络设备按授权范围直接访问远端内存，减少部分软件路径。乐观执行先读取和计算，提交时验证读取版本与资源条件，再保存相应决定。通信更快以后，版本、日志、持久化和故障域共同决定的恢复关系依然要有具体结构。

计算运行时负责把工作放到可用资源上，并管理输入、依赖和结果。MapReduce 将固定输入分成映射、汇聚等阶段，Hadoop 为存储、资源与作业安排不同角色，Spark 用数据集依赖组织计算与重做。推测执行可以让同一逻辑任务出现两个尝试，正式输出由提交与发布协议收束。批量对账又把这些能力用于扫描固定来源，检查付款、订单、事件和履约之间的差异。

Ray 以任务、对象引用和 Actor 组织动态计算。Actor 将实例状态与发往它的消息集中起来，适合复用已加载数据和局部执行状态，具体并发与重启行为由运行时规定。事件函数则由平台按调用或事件取得一段执行时间，AWS Lambda 在执行环境、队列、事件源和函数之间安排接力，Firecracker 等隔离技术承担相应计算边界。长期业务责任从持久来源读取，短生命周期的计算单元便能继续参与恢复。

持久执行进一步保存程序从哪里继续。DBOS 把工作流状态与步骤输出放进数据库，Restate 用独立日志运行时保存调用及其执行历史，LangGraph 用图状态与检查点组织模型、工具和人工等待。Agent 在这里指由模型提出动作、程序调用工具并取得观察的执行系统。模型决定、具体批准与原业务身份先形成可恢复关联，工具重入以后才能接回已经受理的购买。

近期研究还会调整等待与发布的位置。libDSE 让内部推测状态与后台持久化重叠，在对外输出前等待相关依赖可以恢复。Jetpack 在原共识路径之外，为符合条件的命令增加并行快路径。Bodega 把权威读取资格安排到选定地域，相应写入与租约承担新的确认成本。这些方案在自己的故障、顺序和时间条件下改变等待，完整购买继续通过业务关系连接各段决定。

计算与保存分离也会改变存储结构。SlateDB 把日志和有序文件放进对象存储，区分内存可见与持久完成。Iceberg 用快照、清单和目录发布表的文件集合，Delta Lake 用版本化事务日志描述有效文件变化。作业重试时，需要既判断候选文件是否发布，也识别原来源批次是否已经进入目标。它们把逐键状态、整组文件与长期分析数据分别放进相应的恢复路径。

从这些设施回到购买全景，系统设计要回答相同的连接问题：哪个位置拥有事实，哪个位置提供观察，谁保存尚未完成的责任，哪个动作使新状态正式成立。组件的内部结构决定这些接口能够提供的保证与成本。第四部分沿着完整架构深入节点、协议与存储，使这些连接有具体的数据、调用和部署位置。

\section*{把正确性、进展与运行成本放在一起}
\pdfbookmark[1]{把正确性、进展与运行成本放在一起}{overview-operation}

正确性首先需要清楚的业务规格。一次 3000 分购买应当关联相同金额的流水与订单，同一身份恢复原购买，两次合法确认分别进入判断。只有一条扣款记录，却扣成了 4000 分，或者返回了另一个人的订单，同样违反业务含义。安全关系规定在允许的执行历史中，哪些状态和结果应当一直成立。

进展讨论尚未结束的工作怎样最终取得结论。它依赖权威恢复提交能力、网络给出通信机会、旧事务或资格结束占用、工作者得到调度，以及外部平台最终提供能够核验的结果。业务已经提交，客户端获知成功，客户端本地保存了结果，又分别需要对应机会。有限的入口期限和证据保留窗口会改变这些条件。

网络分区会使这些关系直接影响受理。两个位置可以各自在线，却暂时得不到对方的新决定。它们共同支配同一份余额或库存时，新的资源承诺要经过能够仲裁这份资源的保存域。提前分配互不重叠的额度，可以让各地在获配范围内继续受理，相应的委托、使用和回收也要保存。标签等允许合并的状态又可以采取不同安排。业务约束、状态表达和事实归属共同决定分区期间能够继续什么。

分布式系统还要在指定负载下满足时延与吞吐目标。副本之间传输和确认消耗网络与磁盘，跨分区决定增加通信阶段，远端权威读取增加等待。把副本放近读者可以改善部分读取，写入确认与故障切换则会承担相应成本。选择哪种路径，要同时看业务所需观察、访问分布和实际故障域。

容量以工作为单位沿链路计算。每秒的新购买数、同身份重试数、未决查询数和新事件数具有不同含义。平台暂停期间，已受理责任继续增长。平台恢复以后，付款责任下降可能同时带来事件与消费积压上升。每一阶段的完成能力扣除持续流入，才是该阶段可用于排空旧工作的余量。

因此，系统给新受理、原结果查询、确定结果回填、事务恢复和消息交接分别安排预算。准入在建立新持久责任的位置限制新增负担，背压沿积压方向减缓上游，退避与随机等待分散重复尝试。分区与责任类别之间保留调度份额，使热点和长期未知工作都能得到处理机会。

故障隔离也通过这些资源关系成立。支付平台暂停时，查询与既有消息发布可以继续；若所有路径共享一个已经耗尽的连接池，暂停就会扩散。复制、控制、平台和消费路径分别设置可观测的容量与依赖，部署才能把预期中的隔离落实到实际运行。

多地域恢复同时考虑保存点和能力终点。恢复点目标说明允许回到多旧的事实，恢复时间目标说明哪项服务能力何时恢复。原结果可查、新购买开放、未决责任排空与实际仓库交付可以在不同时间完成。恢复预算沿依赖关系安排机器、日志、密钥、网络和外部核对，运维界面分别显示这些进度。

长期运行还包含代码升级、格式兼容和组件维护。新的实现要能解释旧身份、原结果与未完成工作，历史回放按保存的版本继续。软件许可、服务方式和商业支持也进入选型成本。架构选择由此同时涉及业务保证、资源成本与组织能够承担的维护责任。

\section*{用完整历史检验设计}
\pdfbookmark[1]{用完整历史检验设计}{overview-validation}

检验从独立规定允许行为开始。可以先写下一个顺序业务对象：新身份在合法资金和授权条件下成功或拒绝，已有同参身份返回原结论，异参得到冲突，查询报告其观察点的事实。并发实现再与这份规格比较，检查一组交错动作能否解释成符合约定顺序与结果的历史。

线性化提供一种这样的观察要求：每项逻辑操作可以视为在调用与完成之间的某一点生效，已经完成后才开始的另一项操作尊重它们的先后。购买的逻辑区间可以跨越多次通信尝试。提交位置、查询快照和原结果恢复怎样进入这项要求，需要由当前模型与实现逐项说明。跨域工作流的阶段与状态则按它自己的规格解释。

不变量把业务关系写成每个可达状态都应满足的条件。证明先检查初始状态，再说明每项允许迁移怎样保持身份、金额、归属和结果关系。底层约束、原子性与持久性在论证中作为明确前提。加入退款、充值、回收或新的写入路径，迁移集合改变，相应关系就一起扩展。

有界模型可以枚举有限参与者、步骤和状态中的交错，展示旧方案怎样走到错误，也检查修订后相同范围内的行为。实现实验在具体进程、协议和保存位置安排中断，检查实际留下的记录。恢复演练进一步把备份、路由、资格、外部证据和容量带进完整链路。各类方法分别提供对应范围的证据，共同连接业务规格与运行行为。

可观测性让这套连接在上线以后继续使用。从一次请求找到原购买，从原购买找到付款、事件和履约，从未决责任找到最近证据与下一步安排，值守者就能定位等待发生在哪个边界。尝试次数、结果未知的年龄、队列变化和责任完成率各自解释一个方面，观测依据与恢复入口相连，操作才能沿原业务关系推进。

\section*{四部分怎样展开这幅全景}
\pdfbookmark[1]{四部分怎样展开这幅全景}{overview-reading}

第一部分从一个正常运行的下单服务开始，把购买与本地扣款收进同一个事务域。一次提交后的响应丢失，使执行尝试与业务效果分离。随后建立稳定身份、并发仲裁、共同提交与原结果，再用完整历史、不变量和故障实验检验。这里形成全书的基本观察方法：列出实际发生的动作，分别查看各方取得的信息，沿允许的迁移判断下一步。

第二部分把扣款交给外部支付平台，购买开始跨越独立提交边界。付款未知、持久工作流、执行资格、事务内事件与消费者接收，依次把一项工作组织成可以接续的协作。每个接收者在自己的范围内保存效果，发送者通过原身份与确定证据结束相应责任。

第三部分转向这些机制所依赖的记忆。复制让接班副本保存已经确认的决定，回收安排旧请求的准入与剩余证据，状态表达改变合并和历史需求，旧备份与地域切换检验外部历史的延续。退款、跨版本恢复和最终保证又把业务约束、进展条件与验证证据放回同一系统。

第四部分从完整购买架构出发，先确定服务协作、事实所有权和关键链路，再深入消息、分区复制、控制面、缓存、存储、事务与计算运行时。多地域部署、容量、观测和恢复预算使这些关系得到实际承载，工业实现与前沿方案展示保存、执行和等待位置的不同选择。

进入第一章时，用户所见的一次购买已经有了完整背景。后面的每次修改，都可以放回这幅图中：它决定了哪项业务含义，在哪个位置形成事实，沿什么关联继续，又依靠哪些保存与运行条件经受故障。第一章先从其中最小的共同提交范围开始，实现一笔能够正常完成的购买。

\FloatBarrier
\clearpage
\endgroup
